% Material cut from the retrocausal-accordance footnote of golden_reasoning.tex,
% preserved verbatim for the corpus archive. Contains: the teacher quotation,
% the Perimeter/Spekkens passage, the August 21 2022 college-decision passage,
% and the full April 2 2026 journal entry.

Technically speaking though, the ``presentational ability'' trait is what I used to formalize the retrocausal accordance principle in 2025 after my recurring intuition since 2021 popped up in my head unprompted following a presentation I gave for a summer program. The original 2021 version was not actually presentation based (at least solely) but rather was more of a blend of a variety of particularly immense flukes I witnessed throughout childhood, (all conditioning on my having been able to produce the frameworks in the first place, I expect to varying degrees of confidence for each), such as to give an example that I kept having teachers in high school who said ``Ajax is the most brilliant student I have ever taught in 30 years, and it is not close''. So whether or not I actually agree that they are \textit{right} per se, the interesting thing is in the credence warranted by someone even saying such a thing to begin with, in that even such a particular thing being spoken is already a fluke. There are a variety of other flukes I can get into at a later date, some obviously tangled up with the framework and a few others which are just counterfactual-based flukes that I still don't have a particular explanation for, but altogether some of these are immense enough that I certainly expect that they are in some way tangled up with conditioning on whatever my net contribution is toward the Golden Point, as is the idea based on the retrocausal accordance principle. Indeed, in the first few days after I entered college even, I was thinking a lot about how in the counterfactuals where I went to a different school instead, whether each counterfactual would prevent me with avenues toward continuing my research on Nexic reasoning, and in turn influence my capacity to influence the Golden Point. Hence, when I decided whether I was going to make the leap of pursuing Aethic reasoning on August 21, 2022, despite my hedging that ``statistically, it is not likely that I succeed anyway'' \cite{Aethus}---I recognized that I already had a perhaps abnormally fluky vantage with the extrusion principle epiphany, and so if the Golden Point was genuinely conditioned on me finishing the derivation, then suddenly the ``unlikely that I succeed'' statistics is replaced headlong with Golden Optimation, under which it was effectively guaranteed that I would succeed. Hence, why this pushed my credence up just enough that I actually took the shot! And sure enough, I got Aethic reasoning out of it! So there's a good direct practical example of the Route World leap of faith having a direct impact on my decision-making! \par
I should note though that I also worked out some guesses with it on a day in April 2026 which failed (i.e. a blunder on my part), and for an interesting reason: I did all my inference conditioning on a particular thing being important that ended up being trivial, so human error in the Route World leap of faith is arguably just as relevant as before. This also led me to a Platonic-like if-then unpacking of this in June 2026, which I'll post below since it's perhaps helpful as initiating philosophy for the Route World phenomenology:
\begin{quote}
    \textbf{My Koide Formula \cite{koide1983new} Idea WAS ALREADY KNOWN!} \\
\textit{Apr 2, 2026 3:38 AM} \\
Wowzers! Yep, title says it all---it's a rediscovery! Yeah and this also recalibrates quite a few things for me concerning how deep that formula actually was into ontology-no-man's-land, since apparently it was actually perfectly accessible from a modest pure math direction!   \\
But, like, the main thing this really recalibrates for me above all is the philosophical concept of how Golden causality interacts with agency---specifically in the sense that this is almost like the ultimate teaching moment for me that devolving into the anxiety of the choice is, like, almost uncannily unfounded because of the organic feedback of that Golden causality. For example, since about January or February I was considering sending an email out to Rob Spekkens about the Aethic concepts, but my Dad and Juliana specifically told me I shouldn't stir the pot with that kind of strategic unpredictability until I hear back from the master's program at Perimeter. \textit{So}, basically what I recognized from this immediately is that I'm watching Golden causality play out in real time---they themselves are giving me a strategic angle (like physical causality) whereas Golden causality itself, by acting through the medium of that constraint, appears to be firmly telling me ``the time isn't yet right'' (like Optimation causality). So, the sublime paradox of it all is that I now know by reading into the Optimation effects that for some reason I'm going to do something big in the coming few months that conditions on maintaining isolation through it, even though I don't have much of a clue what it will be yet beyond random guessing.   \\
So anyways, this was the exact sort of dilemma I was extrapolating from when deciding whether to post my Koide formula just after midnight yesterday, on April 1, (which itself is absurdly fitting)---basically my anxiety was ``what if whether I post this or not has that same signature of absolute cosmic causal necessity, and yet I have to fall on the right decision here and now without knowing which is right''. What I ended up deciding to do was *not* post it, that is until noon twelve hours later, at which point I called Juliana and asked if she thought posting it would be acceptable, and she said ``yeah it's probably fine''. So, I posted it then. The one Golden causality reframe I noticed already from this is that \textit{even if posting it was the right decision, then not having done it at midnight could not have been an existential catastrophe, because literally I would have that second chance without knowing it anyway, thereby not forcing the whole existential system through the one node there}. BUT, and here's where the recursion gets doubly-potent in this anxiety-liberating sense: \textit{the formula itself was literally discovered already, apparently by a guy named Carl Brannen \cite{Brannen2006TheLM} in 2006}. SO, the correlation at play is that here was this paper which I now learned is effectively peripheral to the Golden causality for this reason, and the only reasoning it wouldn't have been was a construct produced falsely by my own anxious mind! Now I see it is orders of magnitude less causally important than the email events, (apparently, anyway), which retroactively changes the narrative of the anxiety---the Golden causal importance was manufactured by me in the first place, and wasn't actually there.  \\
And then, of course by the way, there's also the Aethic layer to this revelation, which itself is kind of existentially potent lol---\textit{there's the extrusion question now of whether my actually posting the paper is what led to its existential demotion, which in turn led to its having been rediscovered in the past 2006 of my current block universe}. As in, what would my Aethic update look like if I hadn't actually posted the paper at all? Is the Golden causality firm and the mental manufacturing peripheral, and it would have been demoted always (whether or not explicitly discovered earlier, since that's perfectly extrusion-susceptible regardless since it genuinely seems inconsequential (because I rediscovered it anyway, and so had it) as well as originally unknown by me), or could it have actually been more nontrivial to solve in the other reality because I myself wouldn't have provided the exact catalyst to demote it by posting it (that is, because if it is Golden causally important, and like the email I would have posted it too soon, then by the proof by contradiction the correction is preventing its Golden causality importance in the first place, which is exactly what I made considerable by posting it). There's almost a karmic flavor to this as well---I got cocky enough in some sense to post it, so its nontriviality evaporated.   \\
In any case, though, the universe appears to have played the ultimate April Fool's day prank on me! Such is life I guess!
\end{quote}
